% GENERATED FILE. Edit canonical lessons, then run npm run generate:books.
% canonical-source-hash: fnv1a64:bcbd674c9af91f34
% canonical-lessons: ZH-C175-jiaju, ZH-C175-shafa, ZH-C175-mianqian, ZH-W175-mian-face, ZH-C175-guizi, ZH-C175-shujia, ZH-R175-review-170-175

\chapter{Sofa and Shelves}
\label{ch:zh-sofa-and-shelves}
\hlchaptermodality{175}

\begin{chapteropening}
\textbf{By the end of this chapter:} \emph{I can name furniture, a sofa, a cupboard and a bookshelf, and say what is in front of me.}

Complete the last lesson of chapter 175: I can name furniture, a sofa, a cupboard and a bookshelf, and say what is in front of me.
\end{chapteropening}

\section[jiājù]{jiājù — furniture}

\label{lesson:ZH-C175-jiaju}



\begin{quote}
Before the new one: say the Mandarin for a ladder, then the Mandarin for a light bulb.
\end{quote}

\subsection*{You'll want to know: jiājù}
\emph{jiājù} — \textquotedblleft{}furniture\textquotedblright{}.

\begin{quote}
\emph{wǒmen yào mǎi xīn jiājù} — We're going to buy new furniture.
\end{quote}

Its characters come later; for now it is a word for the ear.

\begin{cousinweb}
\emph{jiājù} is built from \zh{家} 'home' and \emph{jù}, both inherited from Middle Chinese, \emph{kae} and \emph{gjuH} (Guangyun, in Baxter's transcription), which glosses the second 'complete' and 'to provide': what a home is provided with.
\end{cousinweb}

\begin{grammarlens}[title={What you've built}]
One more word for everyday things.
\end{grammarlens}

\subsection*{Your turn}
\begin{itemize}
\raggedright
  \item \emph{Say it:} \emph{jiājù}
  \item \emph{Say it:} \emph{jiājù}, once more
  \item \emph{Say it:} \emph{wǒmen yào mǎi xīn jiājù} — We're going to buy new furniture.
  \item \emph{Recall:} say the Mandarin for a ladder, then the Mandarin for a light bulb, then say \emph{jiājù} again
\end{itemize}

\begin{tcolorbox}[breakable,colback=teal!4,colframe=teal!35!black,title={Before you move on}]
What does jiājù mean? (\textquotedblleft{}Furniture\textquotedblright{}.) Say it once more.
\end{tcolorbox}
\section[shāfā]{shāfā — sofa}

\label{lesson:ZH-C175-shafa}



\begin{quote}
Before the new one: say the Mandarin for a light bulb, then the Mandarin for furniture.
\end{quote}

\subsection*{You'll want to know: shāfā}
\emph{shāfā} — \textquotedblleft{}sofa\textquotedblright{}.

\begin{quote}
\emph{wǒmen de shāfā hěn shūfu} — Our sofa is really comfortable.
\end{quote}

Its characters come later; for now it is a word for the ear.

\begin{cousinweb}
\emph{shāfā} is usually traced to English \emph{sofa}: the two characters were chosen for their sound, not for what they mean.
\end{cousinweb}

\begin{grammarlens}[title={What you've built}]
One more word for everyday things.
\end{grammarlens}

\subsection*{Your turn}
\begin{itemize}
\raggedright
  \item \emph{Say it:} \emph{shāfā}
  \item \emph{Say it:} \emph{shāfā}, once more
  \item \emph{Say it:} \emph{wǒmen de shāfā hěn shūfu} — Our sofa is really comfortable.
  \item \emph{Recall:} say the Mandarin for a light bulb, then the Mandarin for furniture, then say \emph{shāfā} again
\end{itemize}

\begin{tcolorbox}[breakable,colback=teal!4,colframe=teal!35!black,title={Before you move on}]
What does shāfā mean? (\textquotedblleft{}Sofa\textquotedblright{}.) Say it once more.
\end{tcolorbox}
\section[\texorpdfstring{\zh{面前} (miànqián)}{面前 (miànqián)}]{\zh{面前} (miànqián) — in front of someone; in someone's presence}

\label{lesson:ZH-C175-mianqian}



\begin{quote}
Before the new one: say the Mandarin for furniture, then the Mandarin for a sofa.
\end{quote}

\subsection*{You'll want to know: \zh{面前}}
\textbf{\zh{面前}} — \emph{miànqián} — \textquotedblleft{}in front of someone; in someone's presence\textquotedblright{}.

\begin{quote}
\textbf{\zh{我面前有一杯茶}} \emph{wǒ miànqián yǒu yì bēi chá} — There's a cup of tea in front of me.
\end{quote}

The new character, in pieces: \zh{面} — nine strokes: a horizontal on top, a short stroke falling to the left, then a frame (a vertical, a horizontal that turns down), two verticals inside it, two short bars between them, and the bottom bar last. The next lesson puts it in your hand.

\begin{cousinweb}
\textbf{\zh{面前}} (\emph{miànqián}) is built from \zh{面} 'face' and \zh{前} 'before', both inherited from Middle Chinese, \emph{mjienH} and \emph{dzen} (Guangyun, in Baxter's transcription), which glosses the first 'to face' and 'the front': before one's face.
\end{cousinweb}

\begin{grammarlens}[title={What you've built}]
One more word for everyday things.
\end{grammarlens}

\subsection*{Your turn}
\begin{itemize}
\raggedright
  \item \emph{Say it:} \emph{miànqián}
  \item \emph{Say it:} \emph{miànqián}, once more
  \item \emph{Say it:} \emph{wǒ miànqián yǒu yì bēi chá} — There's a cup of tea in front of me.
  \item \emph{Recall:} say the Mandarin for furniture, then the Mandarin for a sofa, then say \emph{miànqián} again
\end{itemize}

\begin{tcolorbox}[breakable,colback=teal!4,colframe=teal!35!black,title={Before you move on}]
What does \zh{面前} mean? (\textquotedblleft{}In front of someone; in someone's presence\textquotedblright{}.) Say it once more.
\end{tcolorbox}
\section[\texorpdfstring{\zh{面} (miàn)}{面 (miàn)}]{\zh{面} — a bar, a short stroke, then a frame}

\label{lesson:ZH-W175-mian-face}



\begin{quote}
The word first, then the character inside it.

\begin{itemize}
\raggedright
  \item \emph{Recall:} say \emph{miànqián} — \textbf{R1}, one lesson back, before it decays
\end{itemize}

\begin{quote}
\zh{前} \zh{日}
\end{quote}

\emph{Read it:} each one aloud, with nothing said first — \emph{qián}, \emph{rì}
\end{quote}

\subsection*{Script — a bar, a short stroke, then a frame}
\hlblockfigure{\detokenize{figures/ZH-W175-mian-face-filmstrip.pdf}}{How \zh{面} is written, stroke by stroke}

\begin{quote}
\zh{面}
\end{quote}

\textbf{One piece, nine strokes, built from the outside in.}

\textbf{On top, two strokes.} A horizontal across the top. A short stroke falling to the left under its middle.

\textbf{The frame, two strokes.} A vertical on the left. A horizontal along the top that turns down the right side.

\textbf{The inside, four strokes.} Two verticals. Then two short bars between them, one above the other.

\textbf{Last, the bottom bar} that closes the frame.

2 + 2 + 4 + 1 = nine.

\subsection*{Your turn}
\begin{itemize}
\raggedright
  \item \emph{Write it:} \zh{面}, all nine, with the model in front of you
  \item \emph{Cover:} write it again from memory, and count the pieces, not the strokes
\end{itemize}

\begin{tcolorbox}[breakable,colback=teal!4,colframe=teal!35!black,title={Before you move on}]
How many strokes in \zh{面}? (\textbf{Nine}.) What closes the frame? (\textbf{The bottom bar, last of all}.) What stands inside it? (\textbf{Two verticals, with two short bars between them}.)

Source: \href{https://raw.githubusercontent.com/chanind/hanzi-writer-data/68d10a4b21150cae5e1ebbd223eed289cf32d90c/data/\%E9\%9D\%A2.json}{Hanzi Writer Data \zh{面}, pinned stroke paths and medians}.
\end{tcolorbox}
\section[guìzi]{guìzi — cupboard; cabinet}

\label{lesson:ZH-C175-guizi}



\begin{quote}
Before the new one: say the Mandarin for a sofa, then the Mandarin for in front of someone.
\end{quote}

\subsection*{You'll want to know: guìzi}
\emph{guìzi} — \textquotedblleft{}cupboard; cabinet\textquotedblright{}.

\begin{quote}
\emph{bēizi zài guìzi lǐ} — The cups are in the cupboard.
\end{quote}

Its characters come later; for now it is a word for the ear.

\begin{cousinweb}
\emph{guìzi} is built from \emph{guì} 'a chest' and the noun ending \zh{子}, both inherited from Middle Chinese, \emph{gwijH} and \emph{tsiX} (Guangyun, in Baxter's transcription). In the word the second loses its tone.
\end{cousinweb}

\begin{grammarlens}[title={What you've built}]
One more word for everyday things.
\end{grammarlens}

\subsection*{Your turn}
\begin{itemize}
\raggedright
  \item \emph{Say it:} \emph{guìzi}
  \item \emph{Say it:} \emph{guìzi}, once more
  \item \emph{Say it:} \emph{bēizi zài guìzi lǐ} — The cups are in the cupboard.
  \item \emph{Recall:} say the Mandarin for a sofa, then the Mandarin for in front of someone, then say \emph{guìzi} again
\end{itemize}

\begin{tcolorbox}[breakable,colback=teal!4,colframe=teal!35!black,title={Before you move on}]
What does guìzi mean? (\textquotedblleft{}Cupboard; cabinet\textquotedblright{}.) Say it once more.
\end{tcolorbox}
\section[shūjià]{shūjià — bookshelf}

\label{lesson:ZH-C175-shujia}



\begin{quote}
Before the new one: say the Mandarin for in front of someone, then the Mandarin for a cupboard.
\end{quote}

\subsection*{You'll want to know: shūjià}
\emph{shūjià} — \textquotedblleft{}bookshelf\textquotedblright{}.

\begin{quote}
\emph{zhège shūjià hěn hǎokàn} — This bookshelf looks nice.
\end{quote}

Its characters come later; for now it is a word for the ear.

\begin{cousinweb}
\emph{shūjià} is built from \zh{书} 'book' and \emph{jià} 'frame'. The first is inherited from Middle Chinese \emph{syo} (Guangyun, in Baxter's transcription); the second is listed there as \emph{kaeH}, glossed as the frame of a house: a frame for books.
\end{cousinweb}

\begin{grammarlens}[title={What you've built}]
One more word for everyday things.
\end{grammarlens}

\subsection*{Your turn}
\begin{itemize}
\raggedright
  \item \emph{Say it:} \emph{shūjià}
  \item \emph{Say it:} \emph{shūjià}, once more
  \item \emph{Say it:} \emph{zhège shūjià hěn hǎokàn} — This bookshelf looks nice.
  \item \emph{Recall:} say the Mandarin for in front of someone, then the Mandarin for a cupboard, then say \emph{shūjià} again
\end{itemize}

\begin{tcolorbox}[breakable,colback=teal!4,colframe=teal!35!black,title={Before you move on}]
What does shūjià mean? (\textquotedblleft{}Bookshelf\textquotedblright{}.) Say it once more.
\end{tcolorbox}
\section[(dialogue)]{Looking back: from Pots and Pans to Sofa and Shelves}

\label{lesson:ZH-R175-review-170-175}



\begin{quote}
Say the Mandarin for a cooking pot and a teacup.
\end{quote}

\subsection*{Your turn}
\begin{itemize}
\raggedright
  \item \emph{Say it:} a cooking pot, a teacup, a spoon, a fork, a microwave oven
  \item \emph{Say it:} to switch on, an air conditioner, a remote control, to press a button, to touch
  \item \emph{Say it:} a flavour, sour, bitter, smelling good, bland
  \item \emph{Say it:} to invite, to send regards to, an evening party, a get-together, to pay for everyone
  \item \emph{Say it:} to hang something up, to sort something out, to lift together, a ladder, a light bulb
  \item \emph{Say it:} furniture, a sofa, in front of someone, a cupboard, a bookshelf
  \item \emph{Read it:} \zh{杯} \zh{候} \zh{面} — each one aloud, then a word that holds it
\end{itemize}

\begin{tcolorbox}[breakable,colback=teal!4,colframe=teal!35!black,title={Before you move on}]
A cooking pot? (\emph{guō}.) A teacup? (\textbf{\zh{茶杯}}, \emph{chábēi}.) A spoon? (\emph{sháozi}.) A fork? (\emph{chāzi}.) A microwave oven? (\emph{wēibōlú}.) To switch on? (\textbf{\zh{打开}}, \emph{dǎkāi}.) An air conditioner? (\emph{kōngtiáo}.) A remote control? (\emph{yáokòngqì}.) To press a button? (\emph{àn}.) To touch? (\emph{pèng}.) A flavour? (\emph{wèidao}.) Sour? (\emph{suān}.) Bitter? (\emph{kǔ}.) Smelling good? (\emph{xiāng}.) Bland? (\emph{dàn}.) To invite? (\emph{yāoqǐng}.) To send regards to? (\textbf{\zh{问候}}, \emph{wènhòu}.) An evening party? (\textbf{\zh{晚会}}, \emph{wǎnhuì}.) A get-together? (\emph{jùhuì}.) To pay for everyone? (\textbf{\zh{请客}}, \emph{qǐngkè}.) To hang something up? (\emph{guà}.) To sort something out? (\emph{nòng}.) To lift together? (\emph{tái}.) A ladder? (\emph{tīzi}.) A light bulb? (\emph{dēngpào}.) Furniture? (\emph{jiājù}.) A sofa? (\emph{shāfā}.) In front of someone? (\textbf{\zh{面前}}, \emph{miànqián}.) A cupboard? (\emph{guìzi}.) A bookshelf? (\emph{shūjià}.)
\end{tcolorbox}
